\scanpage{86}
\mspage{19}
\margin{85}
ce qui s'écrit aussi
\[
(87)\qquad
\Pi_P^{D,\tau}\simeq D(\rho)*_{D_1}D(\sigma_a)
\]
\margin{somme amalgamée de groupes sur $D_1=\{1,\widetilde\sigma_a\}$}
où
\[
(88)\qquad
\left\{
\begin{aligned}
D(\rho)&=\underbrace{D^+}_{\substack{\text{engendré}\\\text{par }\rho,\text{ avec}\\\rho^n=1}}\cdot\{1,\widetilde\sigma_a\}
\deleted{\unsure{D(\sigma_a)}}
=\text{groupe diédral d'indice }n\\
&\quad\deleted{\unsure{\simeq}}\qquad
\deleted{\unsure{défini}}\text{ par }\rho\\
D(\sigma_a)&=\{1,\sigma_a\}\times\{1,\widetilde\sigma_a\}
\quad\text{groupe diédral d'indice }2\\
&\hspace{12em}\text{défini par }\sigma_a\\
D_1&=\{1,\widetilde\sigma_a\}
\end{aligned}
\right.
\]
L'écriture (77) s'écrit aussi
\[
(89)\quad
\Pi_P^D\xrightarrow{\sim}
\Pi_{0,3}/\langle\rho_1^2,\rho_\infty^n\rangle
=\Pi_1(\mathbb P^1(\mathbb C),
\underbrace{\infty\cdot(0)+2(1)+n(\infty)}_{\text{signature}};P_+)
\]
en transformant $(\rho,\sigma_{a_0})$ en $(\rho_\infty,\rho_1)$.
Une façon a priori d'obtenir un isomorphisme (89) est de considérer la courbe complexe quotient de $\Sigma_P$ par l'action de $D_P$, \deleted{qui} action \inserted{fidèle,} qui a comme points de ramification les $R_s$ (ramification $2$), les $Q_s$ (ramification $2$), et $P_\omega,P_{\omega'}$ (ramification $n$) -- le quotient est \inserted{une courbe alg. complexe} isomorphe à $\mathbb P^1_{\mathbb C}$, l'isomorphisme étant uniquement déterminé quand on le normalise par
\[
(90)\qquad
\left\{
\begin{aligned}
R_s&\longmapsto0&&(s\in S)\\
Q_a&\longmapsto1&&(a\in A)\\
P_\omega&\longmapsto\infty&&(\omega\in\underline\omega)
\end{aligned}
\right.
\qquad(\Sigma_P\longrightarrow\mathbb P^1_{\mathbb C}).
\]
\unsure{On voit déjà} que la carte envisagée sur $\Sigma_P$ est l'image inverse de la carte universelle, par un morphisme $\Sigma_P\to\mathbb P^1_{\mathbb C}$ bien déterminé, \deleted{qui} [Dans le cas épinglé
\margin{C'est un peu une tricherie, \unsure{car} pour construire cette carte, (\unsure{chose qui n'a pas été faite}), il aurait fallu passer par ce morphisme $\Sigma_P\to\mathbb P^1_{\mathbb C}$, cf plus bas}
